\begin{theorem}[纤维均匀提升是唯一的最大熵截面]\label{thm:conclusion-fold-audit-maxent-section-uniqueness}
固定分辨率 \(m\ge 1\)。对任意宏观分布 \(\nu\in\Delta(X_m)\)，定义其 lift 多面体
\[
\mathcal L_m(\nu):=
\{\mu\in\Delta(\Omega_m):P_m\mu=\nu\}.
\]
则 \(Q_m\nu\) 是 \(\mathcal L_m(\nu)\) 上 Shannon 熵的唯一极大元。更精确地，若记
\[
\Gamma_m(\nu):=\mathbb E_\nu[\log d_m(X)],
\]
则有
\[
H(Q_m\nu)=H(\nu)+\Gamma_m(\nu),
\]
并且对任意 \(\mu\in\mathcal L_m(\nu)\)，
\[
H(Q_m\nu)-H(\mu)
=
\sum_{x\in X_m}\nu(x)\Bigl(\log d_m(x)-H(\mu(\,\cdot\,|x))\Bigr)\ge 0.
\]
等号成立当且仅当 \(\mu(\,\cdot\,|x)\) 在每条纤维 \(\Fold_m^{-1}(x)\) 上均为均匀分布。
\end{theorem}
\begin{proof}
把任意 \(\mu\in\mathcal L_m(\nu)\) 按纤维分解为
\[
\mu(a)=\nu(x)\,\mu_x(a)\qquad (a\in \Fold_m^{-1}(x)),
\]
其中 \(\mu_x\in\Delta(\Fold_m^{-1}(x))\)。由熵的链式分解，
\[
H(\mu)=H(\nu)+\sum_{x\in X_m}\nu(x)H(\mu_x).
\]
另一方面，\(Q_m\nu\) 在每条纤维 \(\Fold_m^{-1}(x)\) 上恰为均匀分布，故
\[
H(Q_m\nu)
=
H(\nu)+\sum_{x\in X_m}\nu(x)\log d_m(x)
=
H(\nu)+\Gamma_m(\nu).
\]
两式相减即得
\[
H(Q_m\nu)-H(\mu)
=
\sum_{x\in X_m}\nu(x)\Bigl(\log d_m(x)-H(\mu_x)\Bigr).
\]
由于有限集上熵在均匀分布处唯一取极大，故对每个 \(x\) 都有
\[
H(\mu_x)\le \log d_m(x),
\]
且等号当且仅当 \(\mu_x\) 为该纤维上的均匀分布。于是结论成立。证毕。
\end{proof}

\begin{theorem}[Fold 审计的精确 KL--Pythagoras 分解]\label{thm:conclusion-fold-audit-kl-pythagoras}
\leanverified{paper\_conclusion\_fold\_audit\_kl\_pythagoras}
对任意 \(\mu\in\Delta(\Omega_m)\) 与任意 \(\nu\in\Delta(X_m)\)，有严格恒等式
\[
D(\mu\|Q_m\nu)
=
D(P_m\mu\|\nu)+D(\mu\|Q_mP_m\mu).
\]
特别地，
\[
D(\mu\|Q_mP_m\mu)=H(Q_mP_m\mu)-H(\mu),
\]
并且
\[
D(\mu\|Q_mP_m\mu)
=
\sum_{x\in X_m}(P_m\mu)(x)\Bigl(\log d_m(x)-H(\mu(\,\cdot\,|x))\Bigr).
\]
\end{theorem}
\begin{proof}
记 \(\mu^\flat:=P_m\mu\in\Delta(X_m)\)，并把 \(\mu\) 的纤维条件分布记为 \(\mu_x:=\mu(\,\cdot\,|x)\)。对任意 \(a\in\Fold_m^{-1}(x)\)，由 \(Q_m\) 的定义有
\[
(Q_m\nu)(a)=\frac{\nu(x)}{d_m(x)},
\qquad
(Q_m\mu^\flat)(a)=\frac{\mu^\flat(x)}{d_m(x)}.
\]
于是
\[
\begin{aligned}
D(\mu\|Q_m\nu)
&=
\sum_{x\in X_m}\sum_{a\in\Fold_m^{-1}(x)}
\mu(a)\log\frac{\mu(a)}{\nu(x)/d_m(x)}\\
&=
\sum_{x\in X_m}\mu^\flat(x)\log\frac{\mu^\flat(x)}{\nu(x)}
+\sum_{x\in X_m}\mu^\flat(x)
\sum_{a\in\Fold_m^{-1}(x)}\mu_x(a)\log\frac{\mu_x(a)}{1/d_m(x)}.
\end{aligned}
\]
第一项正是 \(D(\mu^\flat\|\nu)\)，第二项正是 \(D(\mu\|Q_m\mu^\flat)\)，故
\[
D(\mu\|Q_m\nu)=D(P_m\mu\|\nu)+D(\mu\|Q_mP_m\mu).
\]
取 \(\nu=P_m\mu\) 即得第一条特别公式。

再由上一条定理对 \(\nu=P_m\mu\) 的最大熵截面结论，
\[
H(Q_mP_m\mu)-H(\mu)
=
\sum_{x\in X_m}(P_m\mu)(x)\Bigl(\log d_m(x)-H(\mu_x)\Bigr),
\]
而右端与上式第二项完全相同，故
\[
D(\mu\|Q_mP_m\mu)=H(Q_mP_m\mu)-H(\mu),
\]
并得最后一个显式展开式。证毕。
\end{proof}

\begin{theorem}[\texorpdfstring{$K_m=Q_mP_m$}{Km=QmPm} 是全部审计幂等中的普适 KL 极小元]\label{thm:conclusion-fold-audit-universal-kl-minimal-idempotent}
记
\[
\mathfrak R_m:=
\{\,Q_m\kappa P_m:\ \kappa:\Delta(X_m)\to\Delta(X_m)\ \text{为仿射幂等}\,\}.
\]
则对任意 \(w_m=Q_m\kappa P_m\in\mathfrak R_m\) 与任意 \(\mu\in\Delta(\Omega_m)\)，有
\[
D(\mu\|w_m\mu)
=
D(P_m\mu\|\kappa P_m\mu)+D(\mu\|K_m\mu)
\ge
D(\mu\|K_m\mu).
\]
因而 \(K_m\) 具有如下普适极小性质：它是 \(\mathfrak R_m\) 中唯一一个对所有微观分布 \(\mu\) 都达到最小 KL 损失的幂等元。
\end{theorem}
\begin{proof}
对任意 \(\mu\in\Delta(\Omega_m)\)，令 \(\mu^\flat:=P_m\mu\)。对
\[
w_m\mu=Q_m\kappa\mu^\flat
\]
应用定理 \ref{thm:conclusion-fold-audit-kl-pythagoras}，取 \(\nu=\kappa\mu^\flat\)，便得
\[
D(\mu\|w_m\mu)
=
D(\mu^\flat\|\kappa\mu^\flat)+D(\mu\|Q_m\mu^\flat).
\]
由于 \(Q_m\mu^\flat=Q_mP_m\mu=K_m\mu\)，故
\[
D(\mu\|w_m\mu)
=
D(P_m\mu\|\kappa P_m\mu)+D(\mu\|K_m\mu)
\ge
D(\mu\|K_m\mu).
\]
这证明了普适下界。

若某个 \(w_m=Q_m\kappa P_m\) 对所有 \(\mu\) 都与 \(K_m\) 取得相同损失，则上式中第一项必须对所有 \(\mu\) 恒为零，即
\[
D(P_m\mu\|\kappa P_m\mu)=0\qquad (\forall \mu\in\Delta(\Omega_m)).
\]
而 \(P_m:\Delta(\Omega_m)\twoheadrightarrow \Delta(X_m)\) 为满射，遂等价于
\[
D(\nu\|\kappa\nu)=0\qquad (\forall \nu\in\Delta(X_m)).
\]
因此 \(\kappa\nu=\nu\) 对所有 \(\nu\) 成立，即 \(\kappa=\mathrm{id}_{\Delta(X_m)}\)。于是
\[
w_m=Q_mP_m=K_m.
\]
唯一性得证。证毕。
\end{proof}

\begin{theorem}[除 \texorpdfstring{$K_m$}{Km} 外一切审计幂等都带有点态 KL 见证]\label{thm:conclusion-fold-audit-nonminimal-idempotent-pointwise-kl-witness}
若
\[
w_m=Q_m\kappa P_m\in\mathfrak R_m,
\qquad
w_m\neq K_m,
\]
则存在某个微观分布 \(\mu\in\Delta(\Omega_m)\) 使得
\[
D(\mu\|w_m\mu)>D(\mu\|K_m\mu)=0.
\]
更精确地，取任意满足 \(\kappa\nu\neq \nu\) 的 \(\nu\in\Delta(X_m)\)，并令
\[
\mu:=Q_m\nu,
\]
则有
\[
D(\mu\|w_m\mu)=D(\nu\|\kappa\nu)>0.
\]
\end{theorem}
\begin{proof}
由于 \(w_m\neq K_m\)，故 \(\kappa\neq \mathrm{id}_{\Delta(X_m)}\)，从而存在 \(\nu\in\Delta(X_m)\) 使
\[
\kappa\nu\neq \nu.
\]
取 \(\mu:=Q_m\nu\)。由 \(P_mQ_m=\mathrm{id}_{\Delta(X_m)}\) 得
\[
P_m\mu=\nu,
\qquad
K_m\mu=Q_mP_mQ_m\nu=Q_m\nu=\mu.
\]
因此
\[
D(\mu\|K_m\mu)=D(\mu\|\mu)=0.
\]
另一方面，由定理 \ref{thm:conclusion-fold-audit-universal-kl-minimal-idempotent} 的精确分解，
\[
D(\mu\|w_m\mu)
=
D(P_m\mu\|\kappa P_m\mu)+D(\mu\|K_m\mu)
=
D(\nu\|\kappa\nu).
\]
而 \(\kappa\nu\neq \nu\) 时，相对熵严格正，故
\[
D(\mu\|w_m\mu)=D(\nu\|\kappa\nu)>0.
\]
证毕。
\end{proof}

\begin{theorem}[均匀审计律的行列式熵缺口与算术--几何均值恒等式]\label{thm:conclusion-fold-audit-uniform-law-agm-gap}
记
\[
u_m^X:=\mathrm{Unif}(X_m),\qquad
u_m^\Omega:=\mathrm{Unif}(\Omega_m),\qquad
\widetilde u_m:=Q_m u_m^X.
\]
则有
\[
H(\widetilde u_m)
=
\log|X_m|
+
\frac{1}{|X_m|}\sum_{x\in X_m}\log d_m(x),
\]
从而借助行列式势
\[
Z_m(t):=\det(I+t\mathcal F_m)=\prod_{x\in X_m}(1+t\,d_m(x)),
\]
得到
\[
H(\widetilde u_m)
=
\log|X_m|
+
\frac{1}{|X_m|}
\lim_{t\to\infty}\Bigl(\log Z_m(t)-|X_m|\log t\Bigr).
\]
进一步，
\[
D(\widetilde u_m\|u_m^\Omega)
=
\log\frac{\frac1{|X_m|}\sum_{x\in X_m}d_m(x)}
{\left(\prod_{x\in X_m}d_m(x)\right)^{1/|X_m|}}.
\]
亦即，\(\widetilde u_m\) 到全微观均匀律的相对熵，精确等于纤维多重度谱的算术--几何均值缺口。
\end{theorem}
\begin{proof}
由定理 \ref{thm:conclusion-fold-audit-maxent-section-uniqueness} 对 \(\nu=u_m^X\) 的应用，
\[
H(\widetilde u_m)
=
H(u_m^X)+\mathbb E_{u_m^X}[\log d_m(X)]
=
\log|X_m|+\frac1{|X_m|}\sum_{x\in X_m}\log d_m(x).
\]
又由附录中行列式势的高温端读出公式，
\[
\lim_{t\to\infty}\Bigl(\log Z_m(t)-|X_m|\log t\Bigr)
=
\sum_{x\in X_m}\log d_m(x),
\]
故第二式立即成立。

再因 \(u_m^\Omega(a)=2^{-m}\)，有
\[
D(\widetilde u_m\|u_m^\Omega)=m\log 2-H(\widetilde u_m).
\]
把上一式代入，并用
\[
\sum_{x\in X_m}d_m(x)=|\Omega_m|=2^m
\]
化简，得到
\[
\begin{aligned}
D(\widetilde u_m\|u_m^\Omega)
&=
\log\frac{2^m/|X_m|}
{\exp\bigl(\frac1{|X_m|}\sum_x\log d_m(x)\bigr)}\\
&=
\log\frac{\frac1{|X_m|}\sum_x d_m(x)}
{\left(\prod_x d_m(x)\right)^{1/|X_m|}}.
\end{aligned}
\]
这正是所述算术--几何均值缺口。证毕。
\end{proof}

\begin{theorem}[临界可见证书率是同一行列式势的端点缺口]\label{thm:conclusion-fold-audit-visible-witness-determinant-endpoint-gap}
记 multiplicity determinant
\[
Z_m(t):=\det(I+t\mathcal F_m),
\qquad
\mathcal F_m:=(E_mE_m^\ast)^{-1}.
\]
则临界可见证书率
\[
r_\varphi=\log_2\frac{2}{\varphi}
\]
可严格重写为
\[
r_\varphi=
\lim_{m\to\infty}
\frac1m
\log_2
\frac{Z_m'(0)}
{\displaystyle \lim_{t\to\infty}\frac{\dd}{\dd(\log t)}\log Z_m(t)}.
\]
等价地，
\[
r_\varphi=
\lim_{m\to\infty}
\frac1m\log_2\frac{\operatorname{Tr}(\mathcal F_m)}{\operatorname{rank}(K_m)}.
\]
\end{theorem}
\begin{proof}
由行列式势的两端读出，
\[
Z_m'(0)=\operatorname{Tr}(\mathcal F_m)=\sum_{x\in X_m}d_m(x)=2^m,
\]
且
\[
\lim_{t\to\infty}\frac{\dd}{\dd(\log t)}\log Z_m(t)=|X_m|=F_{m+2}.
\]
故
\[
\frac{Z_m'(0)}
{\displaystyle \lim_{t\to\infty}\frac{\dd}{\dd(\log t)}\log Z_m(t)}
=
\frac{2^m}{F_{m+2}}.
\]
对其取 \(\frac1m\log_2\) 并用 Binet 公式
\[
F_{m+2}=\frac{\varphi^{m+2}-(-\varphi^{-1})^{m+2}}{\sqrt5}
=
\varphi^m\cdot\frac{\varphi^2+o(1)}{\sqrt5},
\]
即得
\[
\lim_{m\to\infty}\frac1m\log_2\frac{2^m}{F_{m+2}}
=
1-\log_2\varphi
=
\log_2\frac2\varphi
=
r_\varphi.
\]

另一方面，由 \(\operatorname{Tr}(\mathcal F_m)=2^m\) 以及纤维反射子 \(K_m\) 的像空间维数等于 \(|X_m|\)，即
\[
\operatorname{rank}(K_m)=|X_m|=F_{m+2},
\]
可得第二式。证毕。
\end{proof}

\begin{theorem}[单一行列式势的三端点刚性]\label{thm:conclusion-single-determinant-potential-three-endpoint-rigidity}
记
\[
Z_m(t):=\det(I+t\mathcal F_m)=\prod_{x\in X_m}(1+t\,d_m(x)).
\]
则同一势函数 \(Z_m\) 同时编码以下三类彼此正交的信息：
\begin{enumerate}
\normalfont
  \item 微观总质量端点：
  \[
  Z_m'(0)=\operatorname{Tr}(\mathcal F_m)=2^m;
  \]
  \item 可见秩端点：
  \[
  \lim_{t\to\infty}\frac{\dd}{\dd(\log t)}\log Z_m(t)=|X_m|=F_{m+2};
  \]
  \item 审计均匀律的熵缺口：
  \[
  D(\widetilde u_m\|u_m^\Omega)
  =
  \log\frac{\frac1{|X_m|}\sum_{x\in X_m}d_m(x)}
  {\left(\prod_{x\in X_m}d_m(x)\right)^{1/|X_m|}}
  =
  \log\frac{\mathrm{AM}(d_m)}{\mathrm{GM}(d_m)}.
  \]
\end{enumerate}
并且临界可见证书率满足
\[
r_\varphi
=
\log_2\frac{2}{\varphi}
=
\lim_{m\to\infty}
\frac1m\log_2
\frac{Z_m'(0)}
{\displaystyle \lim_{t\to\infty}\frac{\dd}{\dd(\log t)}\log Z_m(t)}.
\]
\end{theorem}
\begin{proof}
第一式由定理 \ref{thm:conclusion-fold-audit-visible-witness-determinant-endpoint-gap} 的证明直接给出：
\[
Z_m'(0)=\operatorname{Tr}(\mathcal F_m)=\sum_{x\in X_m}d_m(x)=2^m.
\]
第二式同样来自该定理的高温端读出：
\[
\lim_{t\to\infty}\frac{\dd}{\dd(\log t)}\log Z_m(t)=|X_m|=F_{m+2}.
\]
第三式则是定理 \ref{thm:conclusion-fold-audit-uniform-law-agm-gap} 的结论。最后把前两式代入定理 \ref{thm:conclusion-fold-audit-visible-witness-determinant-endpoint-gap} 的端点缺口公式，即得
\[
r_\varphi
=
\lim_{m\to\infty}
\frac1m\log_2
\frac{Z_m'(0)}
{\displaystyle \lim_{t\to\infty}\frac{\dd}{\dd(\log t)}\log Z_m(t)}.
\]
证毕。
\end{proof}

\begin{theorem}[投影词演算在一次折叠后已饱和]\label{thm:conclusion-fold-audit-onefold-saturation}
\leanverified{paper\_conclusion\_fold\_audit\_onefold\_saturation}
在附录 A.5 的投影词演算中，定义观察侧与分布侧经由审计层因子化的一切幂等分别为
\[
\mathfrak I_m^{\mathrm{obs}}
:=
\{\,\iota_m\beta E_m:\ \beta^2=\beta\in \operatorname{End}(\mathcal B_m)\,\}
\cup
\{\mathrm{id}_{\mathcal A_m}\},
\]
\[
\mathfrak I_m^{\mathrm{mea}}
:=
\{\,Q_m\alpha P_m:\ \alpha^2=\alpha\in \operatorname{End}(\mathbb R^{X_m})\,\}
\cup
\{\mathrm{id}_{\mathbb R^{\Omega_m}}\}.
\]
则对每个非恒等幂等
\[
w_m\in \mathfrak I_m^{\mathrm{obs}}\cup \mathfrak I_m^{\mathrm{mea}},
\]
均有
\[
\operatorname{rank}(w_m)\le |X_m|=F_{m+2},
\]
且等号成立当且仅当
\[
w_m=\Pi_m=\iota_mE_m
\qquad\text{或}\qquad
w_m=K_m=Q_mP_m.
\]
因此，对任意满足 \(\operatorname{rank}(w_m)\ge 1\) 的此类幂等族 \((w_m)\)，必有
\[
\liminf_{m\to\infty}
\frac1m\log_2\frac{2^m}{\operatorname{rank}(w_m)}
\ge r_\varphi.
\]
\end{theorem}
\begin{proof}
由投影词正规形定理，观察侧任意非恒等幂等都可写成
\[
w_m=\iota_m\beta E_m
\qquad (\beta^2=\beta\in\operatorname{End}(\mathcal B_m)).
\]
由于 \(E_m:\mathcal A_m\twoheadrightarrow \mathcal B_m\) 为满射且 \(\iota_m:\mathcal B_m\hookrightarrow \mathcal A_m\) 为单射，故
\[
\operatorname{rank}(w_m)=\operatorname{rank}(\beta)\le \dim\mathcal B_m=|X_m|.
\]
若达到等号，则 \(\beta\) 为 \(\mathcal B_m\) 上满秩幂等，故必为恒等，从而
\[
w_m=\iota_mE_m=\Pi_m.
\]

分布侧同理，任意非恒等幂等都可写成
\[
w_m=Q_m\alpha P_m
\qquad (\alpha^2=\alpha\in\operatorname{End}(\mathbb R^{X_m})).
\]
由 \(P_m:\mathbb R^{\Omega_m}\twoheadrightarrow \mathbb R^{X_m}\) 为满射且 \(Q_m:\mathbb R^{X_m}\hookrightarrow \mathbb R^{\Omega_m}\) 为单射，得
\[
\operatorname{rank}(w_m)=\operatorname{rank}(\alpha)\le |X_m|.
\]
若达到等号，则 \(\alpha\) 为 \(\mathbb R^{X_m}\) 上满秩幂等，故 \(\alpha=\mathrm{id}\)，于是
\[
w_m=Q_mP_m=K_m.
\]

因此对所有非恒等幂等都有
\[
\operatorname{rank}(w_m)\le F_{m+2}.
\]
于是
\[
\frac1m\log_2\frac{2^m}{\operatorname{rank}(w_m)}
\ge
\frac1m\log_2\frac{2^m}{F_{m+2}}.
\]
令 \(m\to\infty\) 并用上一定理，便得
\[
\liminf_{m\to\infty}
\frac1m\log_2\frac{2^m}{\operatorname{rank}(w_m)}
\ge
r_\varphi.
\]
证毕。
\end{proof}

\begin{conjecture}[行列式熵律]\label{conj:conclusion-fold-audit-determinant-entropy-law}
存在常数
\[
\delta_\ast\in\bigl(0,(\log 2)\,r_\varphi\bigr)
\]
使得
\[
D(\widetilde u_m\|u_m^\Omega)
=
\delta_\ast\,m+o(m).
\]
等价地，存在
\[
\gamma_\ast:=(\log 2)\,r_\varphi-\delta_\ast
\]
使
\[
\frac1{|X_m|}\sum_{x\in X_m}\log d_m(x)
=
\gamma_\ast\,m+o(m),
\]
亦即
\[
\lim_{m\to\infty}
\frac{1}{m|X_m|}
\lim_{t\to\infty}
\Bigl(\log Z_m(t)-|X_m|\log t\Bigr)
=
\gamma_\ast.
\]
若该猜想成立，则 multiplicity determinant \(Z_m(t)\) 在高温端读出的平均对数纤维体积具有稳定线性斜率；于是端点缺口 \((\log 2)\,r_\varphi\) 将被严格分裂为
\[
(\log 2)\,r_\varphi
=
\underbrace{\gamma_\ast}_{\text{可由纤维几何保留的平均隐熵}}
+
\underbrace{\delta_\ast}_{\text{任何均匀审计不可消去的线性信息损失}}.
\]
\end{conjecture}
